\documentclass[../../main.tex]{subfiles}%
\graphicspath{{\subfix{./image/}}{\subfix{../../image/}}}

\begin{document}

\section{参考文献与注记}\label{section:ch06-07}

本章的主要依据是 Federer\cite{federer1969}、Liu\cite{liu1977}、Reshetnjak\cite{reshetnjak1968} 与 Stein\cite{stein1970}。关于 $L^p$ 可微性的处理借用了\cite{stein1970} 第 8.1 节的想法。近似可微性在\cite{federer1969} 第 3.1.2--3.1.5 节中讨论。\cref{section:ch06-02}中\cref{theorem:ch06-6-5}的证明由 D. Adams 向我们展示。% XREF p=303 "Section 6.2" | p=303 "Theorem 6.5"

我们对 Aleksandrov 定理（\cref{theorem:ch06-6-9}）的证明遵循\cite{reshetnjak1968}，Whitney 延拓定理（\cref{theorem:ch06-6-10}）则取自\cite{federer1969} 第 3.1.13--3.1.14 节。Lipschitz 连续函数用 $C^1$ 函数逼近一事来自 Simon\cite{simon1984} 第 5.3 节，另见\cite{federer1969} 第 3.1.15 节。Sobolev 函数的逼近则依据 Liu\cite{liu1977}。Fefferman\cite{fefferman2005} 建立了 Whitney 延拓定理的一个精细版本。% XREF p=303 "Aleksandrov's Theorem" | p=303 "Whitney's Extension Theorem"

\end{document}
